\chapter[Ottimizzazione Univariata: Lipschitz, Modelli Locali e Ricerca Dicotomica]{Ottimizzazione Univariata:\\Lipschitzianit\`a, Modelli Locali e Ricerca Dicotomica}
\label{chap:ottimizzazione-univariata-lipschitz-dicotomica}

Il presente capitolo analizza i fondamenti teorici e computazionali dell'ottimizzazione non lineare su funzioni scalari univariata $f: X \to \R$, con $X \subseteq \R$. Partendo dalla formulazione generale del problema di ottimo e dal raccordo concettuale con i modelli quadratici e le tecniche di ricerca lineare (\textit{line search}) propedeutiche all'ottimizzazione multivariata, viene dimostrata l'impossibilit\`a intrinseca di localizzare il minimo globale di funzioni generiche mediante oracoli puntuali, motivando l'introduzione dello ``speed limit'' imposto dalla condizione di Lipschitz-continuit\`a. Si esamina la complessit\`a computazionale sublineare $\BigO{LD/\varepsilon}$ e il principio del \textit{No Free Lunch Theorem} per l'ottimizzazione globale, per poi formalizzare la svolta pragmatica verso l'ottimizzazione locale nei bacini di attrazione unimodali. Attraverso l'uso del modello lineare di primo ordine e delle condizioni di stazionariet\`a, viene infine derivato ed analizzato l'algoritmo di \textbf{Ricerca Dicotomica} (\textit{Dichotomic Search} / bisezione del gradiente), dimostrandone la convergenza lineare, la complessit\`a logaritmica $\BigO{\log(LD/\varepsilon)}$ e la drastica accelerazione rispetto all'esplorazione globale esaustiva.

\FloatBarrier
\section{Introduzione, Contesto e Raccordo Metodologico}
\label{sec:univariata-intro-contesto}

\subsection{Perch\'e Studiare l'Ottimizzazione Univariata?}
L'ottimizzazione di funzioni a una sola variabile scalare ($n = 1$) riveste un ruolo architettonico fondamentale all'interno della matematica computazionale e dell'apprendimento automatico, per tre ragioni metodologiche primarie:
\begin{enumerate}
    \item \textbf{Raccordo con i problemi quadratici e l'algebra lineare:} Nel caso univariato, la minimizzazione di una parabola convessa $f(x) = \frac{1}{2} a x^2 + b x$ (con $a > 0$) \`e computazionalmente banale e si risolve in tempo $\BigO{1}$ annullando la derivata prima ($a x + b = 0 \iff x^* = -b/a$). Come analizzato nei capitoli dedicati all'algebra lineare numerica e all'ottimizzazione multivariata (Capitolo~\ref{chap:opt-multivariata-lineare-quadratica} e Capitolo~\ref{chap:ottimizzazione-quadratica-non-omogenea-gmq}), in dimensione arbitraria $n > 1$ tale problema si generalizza nella risoluzione del sistema lineare simmetrico $Q x = -q$, ponendosi all'esatto crocevia tra algebra lineare e ottimizzazione numerica.
    \item \textbf{Fondamento computazionale per il caso multivariato ($n > 1$):} La quasi totalit\`a dei metodi di discesa per problemi multidimensionali ad alta dimensionalit\`a (come il metodo del gradiente, il gradiente coniugato e i metodi Quasi-Newton) opera selezionando a ogni passo una \textbf{direzione di discesa} ammissibile $d_k \in \R^n$ e riducendo la scelta dell'iterato successivo alla minimizzazione scalare lungo tale semiretta:
    \begin{equation}
        \min_{\alpha > 0} \phi(\alpha) \equiv f(x_k + \alpha d_k).
        \label{eq:line-search-univariate-bridge}
    \end{equation}
    Questa procedura atomica, nota come \textit{line search} (ricerca lineare univariata), \`e a tutti gli effetti un problema di ottimizzazione in una sola variabile.
    \item \textbf{Trattabilit\`a geometrica ed analitica:} Lo spazio $\R$ consente di visualizzare intuitivamente concetti altrimenti astratti in $\R^n$, quali le curvature, i bacini di attrazione, le patologie di non convergenza e i tassi di contrazione degli intervalli di incertezza.
\end{enumerate}

\subsection{Formulazione Generale del Problema}
Un problema di ottimizzazione scalare univariato si formula canonicamente come:
\begin{equation}
    (P) \quad f^* = \min \{ f(x) : x \in X \},
\end{equation}
oppure, sul piano della localizzazione degli argomenti minimizzanti:
\begin{equation}
    x^* \in \argmin_{x \in X} f(x) \equiv \{ x \in X : f(x) \le f(z), \; \forall z \in X \}.
\end{equation}

\subsubsection{Restrizione del Dominio: Il Box Compatto}
Qualora il dominio ammissibile $X$ fosse un insieme arbitrario con cardinalit\`a inaccessibile o $f$ una funzione non computabile nel senso di Turing, il problema $(P)$ risulterebbe insolubile a priori. Per rendere la trattazione formalmente accessibile, si assume che il dominio sia lo spazio euclideo reale $X = \R$ oppure, in uno scenario ancor pi\`u strutturato e favorevole, un intervallo chiuso e limitato (\textbf{compatto}):
\begin{equation}
    X = [x_-, x_+] \subset \R, \quad \text{con } -\infty < x_- < x_+ < +\infty,
\end{equation}
caratterizzato da un diametro finito:
\begin{equation}
    D = x_+ - x_- < \infty.
\end{equation}

\subsubsection{Presenza di un Oracolo Puntuale}
Si assume che la funzione obiettivo sia accessibile mediante un \textbf{oracolo di valutazione puntuale} (\textit{zeroth-order oracle}): dato un qualsiasi punto $x \in X$, l'oracolo restituisce in output il valore numerico scalare $f(x)$.

\begin{noteblock}{Astrazione Computazionale dell'Oracolo vs.\ Realt\`a Ingegneristica}
Nella teoria della complessit\`a algoritmica, l'interrogazione dell'oracolo viene convenzionalmente conteggiata come un'operazione atomica a costo $\BigO{1}$. Nella realt\`a dell'ingegneria computazionale, tuttavia, il costo temporale ed energetico di una singola valutazione pu\`o variare in modo drastico:
\begin{itemize}
    \item pu\`o trattarsi di una semplice espressione algebrica chiusa (eseguita dalla CPU in frazioni di microsecondo);
    \item pu\`o richiedere l'integrazione numerica complessa di equazioni differenziali ordinarie o stocastiche (metodi Monte Carlo);
    \item pu\`o essere il risultato di simulatori fluidodinamici o strutturali su larga scala (\textit{simulation-based optimization}), dove calcolare $f(x)$ richiede ore o giorni di calcolo distribuito su cluster HPC.
\end{itemize}
Di conseguenza, la metrica primaria di efficienza di un algoritmo non \`e unicamente il tempo di CPU accessorio, ma il \textbf{numero totale di chiamate all'oracolo} necessarie a raggiungere un'accuratezza prefissata.
\end{noteblock}

\FloatBarrier
\section{Impossibilit\`a dell'Ottimizzazione Globale per Funzioni Generiche}
\label{sec:univariata-impossibilita-globale}

\subsection{La Patologia dei Minimi Isolati e degli ``Spike'' Arbitrariamente Stretti}
Ci si potrebbe domandare se, restringendosi a un compatto $X = [x_-, x_+]$ di diametro finito $D$ e disponendo di un oracolo puntuale, sia sempre possibile determinare numericamente il minimo globale $f^*$ (anche solo con precisione approssimata). 

La risposta teorica \`e categoricamente \textbf{negativa}:
\begin{itemize}
    \item L'intervallo $[x_-, x_+]$ contiene una quantit\`a non numerabile di punti ($|\R| = \mathfrak{c}$). Qualsiasi algoritmo eseguibile in tempo finito pu\`o interrogare l'oracolo solo su un sottoinsieme finito di campioni $\{x_1, x_2, \dots, x_k\}$.
    \item Se alla funzione non viene imposta alcuna regolarit\`a topologica (funzione discontinua), il valore minimo potrebbe essere concentrato in un singolo punto isolato $x^*$:
    \begin{equation}
        f(x) = \begin{cases} 0 & \text{se } x = x^* \\ 1 & \text{se } x \neq x^* \end{cases}
    \end{equation}
    La probabilit\`a che un campionamento deterministico o stocastico finito intercetti esattamente $x^*$ \`e rigorosamente nulla.
    \item Ancor pi\`u rilevante: \textbf{anche richiedendo la continuit\`a classica ($f \in C^0$)}, una funzione continua pu\`o esibire picchi verso il basso (\textit{spike}) di ampiezza arbitrariamente sottile:
\end{itemize}

\begin{figure}[H]
\centering
\begin{tikzpicture}[scale=1.15, >=Stealth]
    % Griglia leggera di sfondo per riferimento
    \draw[help lines, color=gray!20, dashed] (0, 0) grid (8.5, 3.2);

    % Assi cartesiani
    \draw[->, thick, gray!80!black] (-0.3, 0) -- (8.8, 0) node[below] {$x$};
    \draw[->, thick, gray!80!black] (0, -0.3) -- (0, 3.3) node[left] {$f(x)$};

    % Livello base piatto
    \draw[dashed, gray!60, thick] (0, 2.2) -- (8.4, 2.2);
    \node[left, font=\footnotesize] at (0, 2.2) {$1.0$};
    \node[left, font=\footnotesize] at (0, 0.4) {$0.0$};

    % Estremi intervallo di ricerca compatto [x_-, x_+]
    \draw[thick, gray!90!black] (0.8, -0.15) -- (0.8, 0.15);
    \node[below=4pt, font=\small\bfseries] at (0.8, 0) {$x_-$};
    \draw[thick, gray!90!black] (8.0, -0.15) -- (8.0, 0.15);
    \node[below=4pt, font=\small\bfseries] at (8.0, 0) {$x_+$};

    % Curva f(x) con spike stretto
    \draw[very thick, blue!85!black] (0.8, 2.2) -- (4.25, 2.2) -- (4.45, 0.4) -- (4.65, 2.2) -- (8.0, 2.2);
    \node[above, blue!85!black, font=\small\bfseries] at (2.2, 2.25) {$f(x)$};

    % Punto di minimo globale isolato nello spike
    \filldraw[red!85!black] (4.45, 0.4) circle (2.2pt);
    \draw[dashed, red!75!black, thin] (4.45, 0.4) -- (4.45, 0);
    \node[below=4pt, font=\small\bfseries, red!85!black] at (4.45, 0) {$x^*$};

    % Punti campionati dall'algoritmo (griglia regolare a passo Delta x)
    \foreach \x in {1.4, 2.6, 3.8, 5.0, 6.2, 7.4} {
        \filldraw[teal!80!black] (\x, 2.2) circle (2.2pt);
        \draw[dotted, teal!70!black, thick] (\x, 0) -- (\x, 2.2);
        \filldraw[teal!80!black] (\x, 0) circle (1.5pt);
    }
    \node[below=4pt, font=\scriptsize, teal!80!black] at (2.6, 0) {$x_i$};
    \node[below=4pt, font=\scriptsize, teal!80!black] at (3.8, 0) {$x_{i+1}$};
    \node[below=4pt, font=\scriptsize, teal!80!black] at (5.0, 0) {$x_{i+2}$};

    % Quota ampiezza dello spike (delta)
    \draw[<->, thick, purple!85!black] (4.25, 2.5) -- (4.65, 2.5) node[midway, above=1pt, font=\scriptsize\bfseries] {Ampiezza $\delta \to 0$};
    \draw[dashed, purple!60!black, thin] (4.25, 2.2) -- (4.25, 2.5);
    \draw[dashed, purple!60!black, thin] (4.65, 2.2) -- (4.65, 2.5);

    % Callout esplicativo senza sovrapposizioni
    \node[draw=red!70!black, fill=red!5, rounded corners=3pt, font=\scriptsize, align=left] at (6.6, 1.3) {
        \textbf{Spike non rilevato:}\\
        $\delta < x_{i+2} - x_{i+1}$\\
        $\forall j, \; f(x_j) = 1.0$\\
        L'oracolo nasconde $x^*$!
    };
\end{tikzpicture}
\caption{La patologia dello \textit{spike} continuo: anche su un compatto $[x_-, x_+]$, una funzione continua pu\`o precipitare verso il minimo su una base $\delta$ pi\`u stretta della distanza tra due campionamenti adiacenti, eludendo qualsiasi reticolo finito di nodi.}
\label{fig:pathology-spike}
\end{figure}

Poich\'e la larghezza di base $\delta$ dello spike pu\`o essere resa arbitrariamente pi\`u piccola della distanza minima tra due campioni consecutivi generati da un algoritmo, l'avversario oracolare potr\`a sempre nascondere il minimo tra due punti campionati, restituendo un valore costantemente piatto ($f(x_i) = 1$) su tutti i nodi ispezionati.

\subsection{La Necessit\`a di uno ``Speed Limit''}
L'argomento dello spike dimostra in modo ineccepibile che la sola continuit\`a topologica \`e una propriet\`a troppo debole per fondare l'ottimizzazione numerica. Per rendere il problema computazionalmente trattabile, \`e indispensabile vincolare la rapidit\`a con cui la funzione pu\`o variare: occorre imporre uno \textbf{``speed limit''} (limite superiore invalicabile sul tasso di variazione locale).

Conoscendo il valore della funzione in un punto campionato $x_i$, lo speed limit assicura che in nessun punto adiacente $z$ la funzione possa precipitare al di sotto di una barriera geometrica deterministica, impedendo la formazione di spike arbitrariamente ripidi.

\FloatBarrier
\section{La Condizione di Lipschitz e il Limite di Velocit\`a}
\label{sec:univariata-lipschitz-speed-limit}

\subsection[Definizione di Lipschitz-continuit\`a]{Definizione di Lipschitz-continuit\`a\\(L-continuit\`a)}
La formalizzazione matematica dello speed limit \`e fornita dalla nozione di continuit\`a lipschitziana:

\begin{definition}[Funzione Lipschitz-continua / L-continua]
\label{def:lipschitz-univariata}
Sia $X \subseteq \R$. Una funzione scalare $f: X \to \R$ si definisce \textbf{Lipschitz-continua} (o $L$-continua) su $X$ se esiste una costante reale finita $L > 0$ (detta \textbf{costante di Lipschitz}) tale che:
\begin{equation}
    |f(x) - f(z)| \le L |x - z|, \quad \forall x, z \in X.
    \label{eq:lipschitz-def}
\end{equation}
\end{definition}

\subsubsection{Interpretazione Geometrica}
Dividendo ambo i membri della disequazione~\eqref{eq:lipschitz-def} per $|x - z| > 0$, si ottiene l'equivalente interpretazione geometrica:
\begin{equation}
    \left| \frac{f(x) - f(z)}{x - z} \right| \le L, \quad \forall x \neq z \in X.
\end{equation}
La pendenza (coefficiente angolare) di qualunque retta secante passante per una coppia arbitraria di punti del grafo di $f$ \`e uniformemente limitata in valore assoluto da $L$. Geometricamente, il grafo di $f$ \`e confinato all'interno di un doppio cono con vertice in $(x, f(x))$ e pendenze $\pm L$:
\begin{equation}
    f(x) - L|z - x| \le f(z) \le f(x) + L|z - x|, \quad \forall z \in X.
\end{equation}

\subsubsection{L-continuit\`a Globale vs.\ Locale}
\`E essenziale distinguere l'ambito di validit\`a della propriet\`a:
\begin{itemize}
    \item \textbf{Globalmente $L$-continua:} se la condizione~\eqref{eq:lipschitz-def} sussiste uniformemente su tutto il dominio non limitato $X = \R$ con un'unica costante $L < \infty$.
    \item \textbf{Localmente $L$-continua in $x$:} se per ogni punto $x \in X$ esiste un intorno sferico $B(x, \varepsilon) = [x - \varepsilon, x + \varepsilon]$ tale che la disuguaglianza~\eqref{eq:lipschitz-def} sia soddisfatta per ogni coppia di punti interni a tale intorno, con costante $L(x)$ dipendente dal punto.
\end{itemize}
La costante di Lipschitz $L$ dipende intrinsecamente dall'estensione dell'insieme considerato: all'allargarsi del dominio, il modulo delle pendenze massime ammissibili pu\`o crescere illimitatamente. Dunque, \textbf{essere localmente $L$-continua ovunque non implica affatto essere globalmente $L$-continua}.

\subsection{La L-continuit\`a Implica la Continuit\`a Ordinaria ($C^0$)}
La lipschitzianit\`a \`e una propriet\`a strettamente pi\`u forte della continuit\`a classica secondo Cauchy:

\begin{theorem}[$L\text{-continuit\`a} \implies C^0$]
Se $f: X \to \R$ \`e Lipschitz-continua con costante $L$ su $X$, allora $f$ \`e continua in ogni punto di $X$.
\end{theorem}
\begin{proof}
Sia fissato ad arbitrio un punto $x \in X$ e una tolleranza di errore $\varepsilon > 0$. Poniamo:
\begin{equation}
    \delta = \frac{\varepsilon}{L} > 0.
\end{equation}
Consideriamo un qualsiasi punto $z \in X$ tale che $|z - x| \le \delta$. Applicando la disuguaglianza di Lipschitz~\eqref{eq:lipschitz-def}, si ricava immediatamente:
\begin{equation}
    |f(z) - f(x)| \le L |z - x| \le L \delta = L \left( \frac{\varepsilon}{L} \right) = \varepsilon.
\end{equation}
Dalla definizione di continuit\`a di Cauchy, $f$ \`e continua in $x$. Poich\'e $x$ \`e arbitrario, $f \in C^0(X)$.
\end{proof}

\subsection{Esempi Notevoli e Controesempi Didattici}
Per comprendere i limiti della propriet\`a di Lipschitz, analizziamo due casi didattici emblematici:

\begin{example}[Funzione localmente Lipschitziana ma non globalmente Lipschitziana]
Si consideri la parabola canonica $f(x) = x^2$ su tutto l'asse reale $X = \R$.
\begin{itemize}
    \item Fissato un qualunque compatto $[-M, M]$, il rapporto incrementale soddisfa:
    \begin{equation}
        \left| \frac{f(x) - f(z)}{x - z} \right| = |x + z| \le |x| + |z| \le 2M.
    \end{equation}
    Dunque $f(x)$ \`e localmente $L$-continua su qualsiasi insieme limitato con costante $L = 2M$.
    \item Tuttavia, se consideriamo $X = \R$, la derivata prima $f'(x) = 2x$ diverge per $x \to \infty$. Scegliendo $z = x + 1$, si ha $|f(x+1) - f(x)| = |2x + 1| \to \infty$ per $x \to \infty$. Non esiste alcuno scalare finito $L$ valido uniformemente su tutto $\R$.
\end{itemize}
\end{example}

\begin{example}[Funzione continua ma non Lipschitziana su un insieme compatto]
Si consideri la funzione $f(x) = \sqrt{|x|}$ (oppure $f(x) = x^{2/3}$) sull'intervallo chiuso e limitato $X = [-1, 1]$.
\begin{itemize}
    \item La funzione \`e continua ovunque su $[-1, 1]$ per i teoremi algebrici elementari.
    \item Tuttavia, calcolando il rapporto incrementale nell'intorno dell'origine ponendo $z = 0$ e $x > 0$:
    \begin{equation}
        \frac{|f(x) - f(0)|}{|x - 0|} = \frac{\sqrt{x}}{x} = \frac{1}{\sqrt{x}}.
    \end{equation}
    Per $x \to 0^+$, la quantit\`a $1/\sqrt{x} \to +\infty$. La retta tangente al grafico nell'origine diviene rigorosamente verticale ($f'(0) = \infty$), violando qualsiasi barriera finita $L$.
\end{itemize}
Questo controesempio dimostra che la compattezza del dominio e la continuit\`a non sono condizioni sufficienti a garantire la lipschitzianit\`a: la pendenza locale non deve mai degenerare all'infinito.
\end{example}

\FloatBarrier
\section{Ottimizzazione Globale di Funzioni Lipschitziane}
\label{sec:univariata-ottimizzazione-globale-lipschitz}

\subsection{Algoritmo di Campionamento Uniforme per $\varepsilon$-Ottimalit\`a}
Una volta imposto lo ``speed limit'' $L$ e ristretto il problema a un intervallo compatto $X = [x_-, x_+]$ di diametro $D = x_+ - x_- < \infty$, l'ottimizzazione globale approssimata diviene matematicamente possibile.

\begin{definition}[Punto $\varepsilon$-Ottimo Globale]
Dato $\varepsilon > 0$, un punto ammissibile $\bar{x} \in X$ si definisce \textbf{$\varepsilon$-ottimo globale} per il problema $\min_{x \in X} f(x)$ se:
\begin{equation}
    f(\bar{x}) - f^* \le \varepsilon, \quad \text{dove } f^* = \min_{x \in X} f(x).
\end{equation}
\end{definition}

\subsubsection{Derivazione del Passo di Campionamento $\Delta x$}
Supponiamo di discretizzare l'intervallo $[x_-, x_+]$ mediante un reticolo uniforme di punti equispaziati a distanza costante $\Delta x$:
\begin{equation}
    x_i = x_- + i \cdot \Delta x, \quad i = 0, 1, 2, \dots, k.
\end{equation}

\begin{figure}[H]
\centering
\begin{tikzpicture}[scale=1.0, >=Stealth]
    % Griglia leggera di sfondo
    \draw[help lines, color=gray!20, dashed] (0, 0) grid (9.5, 3.4);

    % Assi coordinati
    \draw[->, thick, gray!80!black] (0, 0) -- (9.8, 0) node[below] {$x$};
    \draw[->, thick, gray!80!black] (0, 0) -- (0, 3.5) node[left] {$f(x)$};

    % Estremi dei campioni consecutivi x_i e x_{i+1}
    \coordinate (XI) at (1.8, 0);
    \coordinate (XIP1) at (7.8, 0);
    \coordinate (XMID) at (4.8, 0);

    % Nodi campionati sull'asse x
    \draw[thick, gray!90!black] (1.8, -0.15) -- (1.8, 0.15);
    \node[below=4pt, font=\small\bfseries, teal!80!black] at (1.8, 0) {$x_i$};
    \draw[thick, gray!90!black] (7.8, -0.15) -- (7.8, 0.15);
    \node[below=4pt, font=\small\bfseries, teal!80!black] at (7.8, 0) {$x_{i+1}$};
    \filldraw[teal!80!black] (1.8, 0) circle (2.2pt);
    \filldraw[teal!80!black] (7.8, 0) circle (2.2pt);

    % Valutazioni sui nodi campionati f(x_i) e f(x_{i+1})
    \coordinate (FXI) at (1.8, 2.5);
    \coordinate (FXIP1) at (7.8, 2.5);
    \draw[dotted, teal!80!black, thick] (XI) -- (FXI);
    \draw[dotted, teal!80!black, thick] (XIP1) -- (FXIP1);
    \filldraw[teal!80!black] (FXI) circle (2.4pt) node[above left, font=\footnotesize\bfseries] {$f(x_i)$};
    \filldraw[teal!80!black] (FXIP1) circle (2.4pt) node[above right, font=\footnotesize\bfseries] {$f(x_{i+1})$};

    % Posizione di x* nel caso peggiore (punto medio esatto)
    \draw[dashed, red!70!black, thin] (XMID) -- (4.8, 0.7);
    \filldraw[red!85!black] (XMID) circle (2.2pt) node[below=4pt, font=\small\bfseries, red!85!black] {$x^*$ (caso peggiore)};
    \filldraw[red!85!black] (4.8, 0.7) circle (2.4pt) node[above=3pt, font=\footnotesize\bfseries, red!85!black] {$f(x^*)$};

    % Cono inferiore Lipschitziano (pendenza massima -L da x_i e +L verso x_{i+1})
    \draw[very thick, orange!90!black] (FXI) -- (4.8, 0.7) node[midway, left=2pt, font=\scriptsize\bfseries] {Pendenza $-L$};
    \draw[very thick, orange!90!black] (4.8, 0.7) -- (FXIP1) node[midway, right=2pt, font=\scriptsize\bfseries] {Pendenza $+L$};

    % Quote orizzontali ben distanziate
    \draw[<->, thick, blue!80!black] (1.8, 3.1) -- (7.8, 3.1) node[midway, above=1pt, font=\small\bfseries] {Passo di campionamento $\Delta x = \frac{2\varepsilon}{L}$};
    \draw[dashed, blue!50, thin] (1.8, 2.5) -- (1.8, 3.1);
    \draw[dashed, blue!50, thin] (7.8, 2.5) -- (7.8, 3.1);

    \draw[<->, thick, red!85!black] (1.8, 0.35) -- (4.8, 0.35) node[midway, above=1pt, font=\scriptsize\bfseries] {$|x^* - x_i| \le \frac{\Delta x}{2}$};

    % Quota verticale dell'errore massimo epsilon
    \draw[<->, thick, purple!85!black] (8.5, 0.7) -- (8.5, 2.5) node[midway, right=2pt, font=\scriptsize\bfseries, align=left] {Errore massimo\\$f(x_i) - f(x^*) \le L \frac{\Delta x}{2} = \varepsilon$};
    \draw[dashed, purple!50, thin] (4.8, 0.7) -- (8.5, 0.7);
    \draw[dashed, purple!50, thin] (7.8, 2.5) -- (8.5, 2.5);
\end{tikzpicture}
\caption{Analisi dell'errore nel caso peggiore per campionamento uniforme: se la funzione \`e $L$-Lipschitziana, nel caso peggiore il minimo $x^*$ si colloca al punto medio $\Delta x / 2$, limitando l'errore di valore a $L \frac{\Delta x}{2} \le \varepsilon$.}
\label{fig:lipschitz-grid-worst-case}
\end{figure}

L'analisi dell'errore nel caso peggiore procede per via deduttiva:
\begin{enumerate}
    \item Il minimo globale incognito $x^*$ cadr\`a necessariamente all'interno di uno dei sottointervalli contigui $[x_i, x_{i+1}]$.
    \item Nel caso peggiore sfavorevole, $x^*$ si colloca esattamente nel punto mediano del sottointervallo. Pertanto, la distanza tra $x^*$ e il punto campionato a esso pi\`u vicino (denotato con $x_{\text{near}}$) \`e maggiorata da:
    \begin{equation}
        |x_{\text{near}} - x^*| \le \frac{\Delta x}{2}.
    \end{equation}
    \item Poich\'e $f$ \`e Lipschitz-continua con costante $L$, la discrepanza tra il valore osservato nel punto campionato e il valore ottimo reale non pu\`o eccedere:
    \begin{equation}
        f(x_{\text{near}}) - f(x^*) \le L |x_{\text{near}} - x^*| \le L \frac{\Delta x}{2}.
    \end{equation}
    \item Imponendo che tale discrepanza non superi la soglia di tolleranza desiderata $\varepsilon$:
    \begin{equation}
        L \frac{\Delta x}{2} \le \varepsilon \iff \Delta x \le \frac{2\varepsilon}{L}.
    \end{equation}
    \item Il massimo passo di campionamento ammissibile che garantisce deterministicamente la $\varepsilon$-ottimalit\`a globale \`e dunque:
    \begin{equation}
        \Delta x = \frac{2\varepsilon}{L}.
    \end{equation}
    \item Il numero complessivo di nodi di campionamento $k$ necessari per coprire l'intero diametro $D = x_+ - x_-$ \`e dato da:
    \begin{equation}
        k = \left\lceil \frac{D}{\Delta x} \right\rceil = \left\lceil \frac{D}{\frac{2\varepsilon}{L}} \right\rceil = \left\lceil \frac{LD}{2\varepsilon} \right\rceil \in \BigO{\frac{LD}{\varepsilon}}.
        \label{eq:complessita-lipschitz-globale}
    \end{equation}
\end{enumerate}

\subsection{Complessit\`a Sublineare e il Teorema del ``No Free Lunch''}

\subsubsection{Complessit\`a Minimax $\Omega(LD/\varepsilon)$}
Il risultato ottenuto non \`e una limitazione specifica dell'algoritmo a griglia uniforme. Nella teoria dell'informazione computazionale~\cite{nesterov2004introductory} \`e rigorosamente dimostrato che, per la classe generale delle funzioni univariate $L$-continue su compatto, \textbf{nessun algoritmo deterministico o stocastico pu\`o scendere al di sotto di $\Omega(LD/\varepsilon)$ valutazioni dell'oracolo nel caso peggiore}.

\subsubsection{Convergenza Sublineare: Il Costo Decimal-by-Decimal}
La complessit\`a~\eqref{eq:complessita-lipschitz-globale} cresce in modo inversamente proporzionale all'errore $\varepsilon$ ($k \propto 1/\varepsilon$):
\begin{itemize}
    \item Se per ottenere un'accuratezza modesta $\varepsilon = 10^{-3}$ occorrono $k = 100$ iterazioni;
    \item per guadagnare una singola cifra decimale ($\varepsilon = 10^{-4}$) occorrono $k = 1\,000$ iterazioni;
    \item per $\varepsilon = 10^{-5}$ ne servono $10\,000$;
    \item per raggiungere la precisione standard dei calcolatori scientifici a 32 bit ($\varepsilon = 10^{-7}$), occorrono oltre \textbf{$10$ milioni ($10^7$) di valutazioni}.
\end{itemize}
Ogni singola cifra decimale di precisione aggiuntiva moltiplica il numero di chiamate oracolari per un fattore $10$. Questo regime di convergenza \`e definito \textbf{sublineare}.

\subsubsection{La Maledizione della Dimensione ($n > 1$)}
Se questo scenario appare gi\`a severo per $n = 1$, nello spazio multidimensionale $\R^n$ il costo computazionale subisce un collasso esponenziale: coprire un ipercubo di diametro $D$ richiede di campionare un reticolo cartesiano contenente:
\begin{equation}
    k_n = \left( \frac{LD}{2\varepsilon} \right)^n \in \BigO{\left(\frac{1}{\varepsilon}\right)^n} \text{ valutazioni}.
\end{equation}
Per $n = 10$ ed $\varepsilon = 10^{-2}$, il numero di interrogazioni eccede $10^{20}$, rendendo il campionamento globale totalmente impraticabile.

\subsubsection{No Free Lunch Theorem (NFL)}
Il docente rimarca che, in virt\`u del celebre \textbf{No Free Lunch Theorem}, non esiste alcun algoritmo universale in grado di superare questo limite: qualsiasi euristica progettata per esplorare pi\`u rapidamente determinate classi di funzioni Lipschitziane risulter\`a sistematicamente pi\`u lenta su altre configurazioni avversarie.

\subsubsection{La Metafora Didattica del ``Bombardamento a Tappeto''}
Per illustrare plasticamente l'inefficienza del campionamento uniforme, il docente ricorre a una metafora bellica: non avendo alcuna informazione preliminare su dove si trovi l'obiettivo nemico all'interno di un'area di un chilometro quadrato, l'unica strategia sicura per colpirlo consiste nell'eseguire un massiccio bombardamento a tappeto, sparando proiettili a reticolo su ogni singolo centimetro quadro di terreno. Si tratta di un approccio a pura forza bruta: esaustivo e sicuro, ma computazionalmente insostenibile.

\FloatBarrier
\section{Il Compromesso dell'Ottimizzazione Locale e i Bacini di Attrazione}
\label{sec:univariata-compromesso-locale}

\subsection{Il ``Trucco alla Captain Kirk''}
Di fronte a una barriera teorica insormontabile come l'intrattabilit\`a dell'ottimizzazione globale per classi generali di funzioni, la matematica computazionale adotta la celebre strategia attribuita al Capitano James T. Kirk in \textit{Star Trek} di fronte alla simulazione impossibile del \textit{Kobayashi Maru}: \textbf{cambiare le regole del gioco}.

Poich\'e localizzare il minimo globale richiede un numero inaccettabile di valutazioni, si rinuncia all'ottimalit\`a globale per convergere efficientemente a un \textbf{minimo locale}.

\subsection{Definizioni Fondamentali: Minimi Locali e Unimodalit\`a}

\begin{definition}[Minimo Locale e Minimo Locale Stretto]
Sia $f: X \to \R$ e sia $x_* \in X$.
\begin{enumerate}
    \item $x_*$ \`e un \textbf{minimo locale} per $f$ se esiste un raggio $\varepsilon > 0$ tale che:
    \begin{equation}
        f(x_*) \le f(x), \quad \forall x \in X \cap [x_* - \varepsilon, x_* + \varepsilon].
    \end{equation}
    \item $x_*$ \`e un \textbf{minimo locale stretto} se vale la disuguaglianza stretta per ogni punto distinto:
    \begin{equation}
        f(x_*) < f(x), \quad \forall x \in X \cap [x_* - \varepsilon, x_* + \varepsilon] \setminus \{x_*\}.
    \end{equation}
\end{enumerate}
\end{definition}

\begin{definition}[Funzione Unimodale]
Una funzione continua $f: [x_-, x_+] \to \R$ si definisce \textbf{unimodale} sull'intervallo se ammette un unico punto di minimo $x_* \in [x_-, x_+]$ ed \`e:
\begin{itemize}
    \item strettamente decrescente alla sinistra dell'ottimo, per ogni $x \in [x_-, x_*]$;
    \item strettamente crescente alla destra dell'ottimo, per ogni $x \in [x_*, x_+]$.
\end{itemize}
\end{definition}

\subsection{Bacini di Attrazione (Attraction Basins)}
Le funzioni obiettivo reali presentano tipicamente una topologia complessa con molteplici oscillazioni, minimi locali, massimi e punti di flesso. Tuttavia, se restringiamo il raggio d'azione all'intorno di uno specifico minimo locale, la funzione assume una configurazione geometrica regolare e prevedibile:

\begin{definition}[Bacino di Attrazione]
Dato un minimo locale $x_*$, il suo \textbf{bacino di attrazione} \`e il massimo sottointervallo aperto $B(x_*) \subseteq X$ contenente $x_*$ all'interno del quale la funzione risulta unimodale.
\end{definition}

\begin{figure}[H]
\centering
\begin{tikzpicture}[scale=1.1, >=Stealth]
    % Griglia leggera di sfondo
    \draw[help lines, color=gray!15, dashed] (0, 0) grid (10.0, 3.8);

    % Assi coordinati
    \draw[->, thick, gray!80!black] (0, 0) -- (10.4, 0) node[below] {$x$};
    \draw[->, thick, gray!80!black] (0, 0) -- (0, 4.0) node[left] {$f(x)$};

    % Regione del bacino di attrazione ombreggiata (da cresta x=3.0 a cresta x=7.0)
    \fill[purple!8] (3.0, 0) -- (3.0, 2.7) 
        .. controls (3.8, 2.7) and (4.2, 0.5) .. (5.0, 0.5)
        .. controls (5.8, 0.5) and (6.2, 2.7) .. (7.0, 2.7)
        -- (7.0, 0) -- cycle;

    % Curva multi-modale f(x) perfettamente liscia (spline di Bézier)
    \draw[very thick, blue!85!black] 
        (0.6, 3.2) .. controls (1.1, 2.0) and (1.3, 1.3) .. (1.8, 1.3)
        .. controls (2.3, 1.3) and (2.5, 2.7) .. (3.0, 2.7)
        .. controls (3.8, 2.7) and (4.2, 0.5) .. (5.0, 0.5)
        .. controls (5.8, 0.5) and (6.2, 2.7) .. (7.0, 2.7)
        .. controls (7.5, 2.7) and (7.7, 1.3) .. (8.2, 1.3)
        .. controls (8.7, 1.3) and (8.9, 2.0) .. (9.4, 3.2);
    \node[above right, blue!85!black, font=\bfseries] at (9.4, 3.2) {$f(x)$};

    % Separatori di bacino (creste/massimi locali) - Posizionati SOPRA il piano cartesiano
    \node (crest1) [draw=purple!70!black, fill=purple!5, rounded corners=3pt, font=\scriptsize, align=center, inner sep=3pt] at (3.0, 4.65) {
        \textbf{Massimo locale}\\(Cresta di separazione)
    };
    \node (crest2) [draw=purple!70!black, fill=purple!5, rounded corners=3pt, font=\scriptsize, align=center, inner sep=3pt] at (7.0, 4.65) {
        \textbf{Massimo locale}\\(Cresta di separazione)
    };
    \draw[dashed, purple!80!black, thick] (3.0, 0) -- (crest1.south);
    \draw[dashed, purple!80!black, thick] (7.0, 0) -- (crest2.south);
    \filldraw[purple!80!black] (3.0, 2.7) circle (2.2pt);
    \filldraw[purple!80!black] (7.0, 2.7) circle (2.2pt);

    % Quota Bacino di Attrazione in alto
    \draw[<->, very thick, purple!85!black] (3.0, 3.35) -- (7.0, 3.35);
    \node[above=2pt, font=\scriptsize\bfseries, align=center, purple!85!black] at (5.0, 3.35) {
        Bacino di Attrazione\\$B(x^*_{\text{glob}})$
    };

    % Minimo locale 1
    \filldraw[teal!80!black] (1.8, 1.3) circle (2.4pt);
    \node (minloc1) [draw=teal!70!black, fill=teal!5, rounded corners=3pt, font=\scriptsize, align=center, inner sep=3.5pt] at (1.8, -0.75) {
        \textbf{Minimo locale}\\$x = x_{\text{loc}, 1}$
    };
    \draw[dashed, teal!70!black, thin] (1.8, 1.3) -- (minloc1.north);

    % Minimo globale
    \filldraw[red!85!black] (5.0, 0.5) circle (2.6pt);
    \node (minglob) [draw=red!75!black, fill=red!5, rounded corners=3pt, font=\scriptsize, align=center, inner sep=3.5pt] at (5.0, -0.75) {
        \textbf{Minimo globale}\\$x = x^*_{\text{glob}}$
    };
    \draw[dashed, red!75!black, thick] (5.0, 0.5) -- (minglob.north);

    % Minimo locale 2
    \filldraw[teal!80!black] (8.2, 1.3) circle (2.4pt);
    \node (minloc2) [draw=teal!70!black, fill=teal!5, rounded corners=3pt, font=\scriptsize, align=center, inner sep=3.5pt] at (8.2, -0.75) {
        \textbf{Minimo locale}\\$x = x_{\text{loc}, 2}$
    };
    \draw[dashed, teal!70!black, thin] (8.2, 1.3) -- (minloc2.north);
    
    % Box didattico posizionato SOTTO il piano cartesiano con ampio margine verticale
    \node (unimodalbox) [draw=purple!70!black, fill=purple!5, rounded corners=4pt, font=\footnotesize, align=center, inner sep=6pt] at (5.0, -2.25) {
        \textbf{Propriet\`a della Regione Unimodale nel Bacino $B(x^*_{\text{glob}})$:}\\[4pt]
        $\begin{aligned}
            x < x^*_{\text{glob}} &\implies f'(x) < 0 \quad \text{(discesa verso destra)} \\
            x > x^*_{\text{glob}} &\implies f'(x) > 0 \quad \text{(discesa verso sinistra)}
        \end{aligned}$
    };

    % Margine di respiro inferiore
    \path (0, -3.15);
\end{tikzpicture}
\caption{Paesaggio non convesso con molteplici minimi locali: all'interno del bacino di attrazione $B(x^*_{\text{glob}})$ delimitato dalle creste adiacenti la funzione \`e unimodale, garantendo la convergenza deterministica dei metodi locali.}
\label{fig:attraction-basins}
\end{figure}

\subsubsection{Il Limite Strutturale dell'Ottimizzazione Locale}
Tutti i minimi locali si presentano localmente indistinguibili: la funzione scende man mano che ci si avvicina al punto di minimo da ambo i lati. Un algoritmo basato su informazioni locali che entra in un determinato bacino converger\`a inesorabilmente al fondo di quel bacino, senza possedere alcuna capacit\`a analitica di stabilire se quel minimo locale coincida con l'ottimo globale o se esista altrove una valle ancor pi\`u profonda.

\subsubsection{Connessione con il Deep Learning Moderno}
Il professore evidenzia come questo compromesso sia esattamente il fondamento operativo dell'addestramento dei modelli di intelligenza artificiale moderna (reti neurali profonde e LLM): la funzione di perdita empirica su miliardi di parametri possiede un numero sterminato di minimi locali non equivalenti. Gli algoritmi di ottimizzazione del primo ordine (come SGD e Adam) si accontentano di convergere a un minimo locale, il quale, per evidenze empiriche e propriet\`a di sovra-parametrizzazione, garantisce performance eccellenti sul compito da apprendere.

\FloatBarrier
\section{Il Modello di Primo Ordine, Derivate e Punti Stazionari}
\label{sec:univariata-modello-primo-ordine-derivate}

\subsection{Usare la Derivata come Guida di Navigazione}
All'interno di un bacino unimodale, per raggiungere il minimo locale \`e sufficiente muoversi seguendo la pendenza locale.
\begin{itemize}
    \item Per una funzione lineare $f(x) = bx + c$, il coefficiente angolare $b$ determina in modo univoco la direzione: se $b > 0$, la retta scende verso sinistra; se $b < 0$, scende verso destra.
    \item Per una funzione non lineare differenziabile, il \textbf{modello lineare di primo ordine} attorno a un punto $x$ \`e rappresentato dalla retta tangente:
    \begin{equation}
        L_x(z) = f(x) + f'(x)(z - x).
    \end{equation}
\end{itemize}
La derivata prima $f'(x)$ funge da bussola direzionale:
\begin{itemize}
    \item $f'(x) < 0 \implies$ la funzione \`e localmente decrescente: il minimo locale si trova a destra ($z > x$);
    \item $f'(x) > 0 \implies$ la funzione \`e localmente crescente: il minimo locale si trova a sinistra ($z < x$);
    \item $f'(x) = 0 \implies$ la retta tangente \`e orizzontale: punto stazionario candidato.
\end{itemize}

\subsection{Differenziabilit\`a, Continuit\`a e Calcolo delle Derivate}
Ricordiamo la definizione formale di derivata:
\begin{equation}
    f'(x) = \lim_{t \to 0} \frac{f(x+t) - f(x)}{t}.
\end{equation}
Come noto dall'analisi matematica elementare, la differenziabilit\`a implica la continuit\`a ordinaria:
\begin{equation}
    \lim_{t \to 0} \big( f(x+t) - f(x) \big) = \lim_{t \to 0} \left( t \cdot \frac{f(x+t) - f(x)}{t} \right) = 0 \cdot f'(x) = 0.
\end{equation}

\subsubsection{Disponibilit\`a Pratica delle Derivate nel Software}
Per funzioni costituite da combinazioni di operatori elementari (algebrici, trigonometrici, esponenziali), la derivata prima \`e ottenibile in forma analitica esatta oppure calcolata in modo ultra-efficiente mediante strumenti software di \textbf{differenziazione automatica} (\textit{AutoDiff} in PyTorch, JAX o comandi simbolici quali \texttt{diff} in MATLAB/Octave).

\subsection{Condizioni di Ottimalit\`a del Primo e del Secondo Ordine}

\begin{theorem}[Condizione Necessaria del Primo Ordine]
Sia $f: X \to \R$ derivabile in un intorno del punto interno $x_* \in \operatorname{int}(X)$. Se $x_*$ \`e un minimo locale per $f$, allora necessariamente:
\begin{equation}
    f'(x_*) = 0.
\end{equation}
\end{theorem}

L'annullamento della derivata prima individua un \textbf{punto stazionario}, ma costituisce una condizione puramente necessaria e non sufficiente: $f'(x_*) = 0$ pu\`o caratterizzare un minimo locale, un massimo locale o un punto di flesso a tangente orizzontale (punto di sella).

\subsubsection{Discriminazione Mediante la Derivata Seconda $f''(x_*)$}
Se la funzione \`e due volte differenziabile ($f \in C^2$), la concavit\`a locale discrimina la natura del punto stazionario:
\begin{itemize}
    \item \textbf{Minimo locale stretto:} $f'(x_*) = 0$ e $f''(x_*) > 0$. La derivata prima sta crescendo (passa da valori negativi a sinistra a valori positivi a destra: concavit\`a verso l'alto).
    \item \textbf{Massimo locale stretto:} $f'(x_*) = 0$ e $f''(x_*) < 0$. La derivata prima sta decrescendo (passa da valori positivi a negativi: concavit\`a verso il basso).
    \item \textbf{Test inconclusivo / Sella:} $f'(x_*) = 0$ e $f''(x_*) = 0$ (es. $f(x) = x^3$ ha un flesso in $x=0$; $f(x) = x^4$ ha un minimo).
\end{itemize}

\begin{figure}[H]
\centering
\begin{tikzpicture}[scale=1.05, >=Stealth]
    % =========================================================================
    % RIGA 1: MINIMO LOCALE STRETTO
    % =========================================================================
    \begin{scope}[yshift=0cm]
        % Box intestazione (y=3.4 per ampio stacco dall'asse y)
        \node[draw=blue!70!black, fill=blue!5, rounded corners=4pt, font=\small, align=center, inner sep=4pt] at (3.0, 3.4) {
            \textbf{Minimo Locale Stretto: Condizioni di Ottimalit\`a}\\
            $f'(x_*) = 0 \quad \text{e} \quad f''(x_*) > 0$
        };

        % Assi coordinati (asse y alto fino a 2.2 con etichetta a sinistra)
        \draw[->, thick, gray!75!black] (-0.2, 0) -- (6.5, 0) node[below] {$x$};
        \draw[->, thick, gray!75!black] (0, -0.2) -- (0, 2.2);
        \node[left=2pt, font=\footnotesize, gray!80!black] at (0, 2.2) {$f(x)$};

        % Parabola convessa (concavità verso l'alto)
        \draw[very thick, blue!85!black, smooth, domain=1.2:4.8] plot (\x, {0.35*(\x-3.0)*(\x-3.0) + 0.6});

        % Punto di minimo x*
        \filldraw[red!85!black] (3.0, 0.6) circle (2.5pt);
        \draw[dashed, red!70!black, thick] (3.0, 0.6) -- (3.0, 0);
        \node[below=2pt, font=\small\bfseries, red!85!black] at (3.0, 0) {$x_*$};
        \node[above=3pt, font=\scriptsize\bfseries, red!85!black, fill=white, fill opacity=0.9, text opacity=1, inner sep=1.5pt, rounded corners=2pt] at (3.0, 0.6) {$(x_*, f(x_*))$};

        % Tangente orizzontale a f(x*)
        \draw[dashed, red!85!black, very thick] (1.4, 0.6) -- (4.6, 0.6);
        \node[right=3pt, font=\scriptsize\bfseries, red!85!black] at (4.6, 0.6) {Tangente: $f'(x_*) = 0$};

        % Pendenza a sinistra: f'(x) < 0
        \draw[->, line width=1.3pt, orange!90!black] (1.5, 1.6) -- (2.2, 0.95);
        \node[above left=1pt, font=\scriptsize\bfseries, orange!90!black, align=center, fill=white, fill opacity=0.9, text opacity=1, inner sep=1.5pt, rounded corners=2pt] at (1.7, 1.5) {Decrescente\\$f'(x) < 0$};

        % Pendenza a destra: f'(x) > 0
        \draw[->, line width=1.3pt, teal!85!black] (3.8, 0.95) -- (4.5, 1.6);
        \node[above right=1pt, font=\scriptsize\bfseries, teal!85!black, align=center, fill=white, fill opacity=0.9, text opacity=1, inner sep=1.5pt, rounded corners=2pt] at (4.3, 1.5) {Crescente\\$f'(x) > 0$};

        % Sintesi concavità
        \node[font=\footnotesize, align=center, text=gray!90!black, text width=12.2cm] at (3.0, -0.85) {
            \textbf{Concavit\`a verso l'alto ($f''(x_*) > 0$):} derivata prima strettamente crescente (da negativa a positiva).
        };
    \end{scope}

    % =========================================================================
    % RIGA 2: MASSIMO LOCALE STRETTO
    % =========================================================================
    \begin{scope}[yshift=-5.6cm]
        % Box intestazione
        \node[draw=teal!70!black, fill=teal!5, rounded corners=4pt, font=\small, align=center, inner sep=4pt] at (3.0, 3.4) {
            \textbf{Massimo Locale Stretto: Condizioni di Ottimalit\`a}\\
            $f'(x_*) = 0 \quad \text{e} \quad f''(x_*) < 0$
        };

        % Assi coordinati
        \draw[->, thick, gray!75!black] (-0.2, 0) -- (6.5, 0) node[below] {$x$};
        \draw[->, thick, gray!75!black] (0, -0.2) -- (0, 2.2);
        \node[left=2pt, font=\footnotesize, gray!80!black] at (0, 2.2) {$f(x)$};

        % Parabola concava (concavità verso il basso)
        \draw[very thick, teal!85!black, smooth, domain=1.2:4.8] plot (\x, {-0.35*(\x-3.0)*(\x-3.0) + 2.1});

        % Punto di massimo x*
        \filldraw[red!85!black] (3.0, 2.1) circle (2.5pt);
        \draw[dashed, red!70!black, thick] (3.0, 2.1) -- (3.0, 0);
        \node[below=2pt, font=\small\bfseries, red!85!black] at (3.0, 0) {$x_*$};
        \node[above=3pt, font=\scriptsize\bfseries, red!85!black, fill=white, fill opacity=0.9, text opacity=1, inner sep=1.5pt, rounded corners=2pt] at (3.0, 2.1) {$(x_*, f(x_*))$};

        % Tangente orizzontale a f(x*)
        \draw[dashed, red!85!black, very thick] (1.4, 2.1) -- (4.6, 2.1);
        \node[right=3pt, font=\scriptsize\bfseries, red!85!black] at (4.6, 2.1) {Tangente: $f'(x_*) = 0$};

        % Pendenza a sinistra: f'(x) > 0
        \draw[->, line width=1.3pt, teal!85!black] (1.5, 0.95) -- (2.2, 1.6);
        \node[below left=1pt, font=\scriptsize\bfseries, teal!85!black, align=center, fill=white, fill opacity=0.9, text opacity=1, inner sep=1.5pt, rounded corners=2pt] at (1.7, 1.1) {Crescente\\$f'(x) > 0$};

        % Pendenza a destra: f'(x) < 0
        \draw[->, line width=1.3pt, orange!90!black] (3.8, 1.6) -- (4.5, 0.95);
        \node[below right=1pt, font=\scriptsize\bfseries, orange!90!black, align=center, fill=white, fill opacity=0.9, text opacity=1, inner sep=1.5pt, rounded corners=2pt] at (4.3, 1.1) {Decrescente\\$f'(x) < 0$};

        % Sintesi concavità
        \node[font=\footnotesize, align=center, text=gray!90!black, text width=12.2cm] at (3.0, -0.85) {
            \textbf{Concavit\`a verso il basso ($f''(x_*) < 0$):} derivata prima strettamente decrescente (da positiva a negativa).
        };
    \end{scope}
\end{tikzpicture}
\caption{Confronto geometrico tra minimo e massimo locale sulla base del segno della derivata seconda $f''(x_*)$: la concavit\`a locale determina la transizione di segno della derivata prima nell'intorno del punto stazionario.}
\label{fig:first-second-order-conditions}
\end{figure}

A livello computazionale, per individuare un minimo locale \`e sufficiente identificare un intervallo in cui la derivata prima transiti da valori negativi a valori positivi, riducendo l'ottimizzazione alla ricerca dello zero di $f'(x)$.

\FloatBarrier
\section{L'Algoritmo di Ricerca Dicotomica (Dichotomic Search)}
\label{sec:univariata-ricerca-dicotomica}

\subsection{Fondamento Teorico: Teorema dei Valori Intermedi}
L'algoritmo di \textbf{Ricerca Dicotomica} (\textit{Dichotomic Search}, o metodo di bisezione applicato alla derivata) poggia su un solido teorema dell'analisi classica:

\begin{theorem}[Teorema di Bolzano / Valori Intermedi applicato alla Derivata]
Sia $f: [x_-, x_+] \to \R$ una funzione derivabile con derivata prima $f'$ continua ($f \in C^1$). Se agli estremi dell'intervallo sussiste la condizione di segno discorde:
\begin{equation}
    f'(x_-) < 0 \quad \text{e} \quad f'(x_+) > 0,
    \label{eq:condizione-bracketing-derivata}
\end{equation}
allora esiste almeno un punto interno $x_* \in (x_-, x_+)$ in cui:
\begin{equation}
    f'(x_*) = 0.
\end{equation}
Inoltre, poich\'e la funzione entra nell'intervallo decrescendo a sinistra ed esce crescendo a destra, l'intervallo $[x_-, x_+]$ racchiude necessariamente almeno un \textbf{minimo locale} di $f$.
\end{theorem}

\subsection{Funzionamento e Pseudocodice dell'Algoritmo}
L'algoritmo opera mantenendo a ogni iterazione l'invariante di segno agli estremi dell'intervallo corrente $[x_-, x_+]$:
\begin{equation}
    \text{Invariante: } f'(x_-) \le -\varepsilon \quad \text{e} \quad f'(x_+) \ge \varepsilon.
\end{equation}
A ogni passo, l'algoritmo valuta la derivata nel punto medio $x = \frac{x_- + x_+}{2}$ e dimezza deterministicamente l'intervallo:
\begin{itemize}
    \item Se $|f'(x)| \le \varepsilon$, il punto medio $x$ \`e un punto stazionario $\varepsilon$-approssimato e la procedura termina;
    \item Se $f'(x) < -\varepsilon$, il minimo locale giace strettamente a destra: si scarta la met\`a sinistra ponendo $x_- \gets x$;
    \item Se $f'(x) > \varepsilon$, il minimo locale giace strettamente a sinistra: si scarta la met\`a destra ponendo $x_+ \gets x$.
\end{itemize}

\begin{algorithm}[H]
\caption{Ricerca Dicotomica per Ottimizzazione Univariata (\texttt{DichotomicSearch})}
\label{alg:ricerca-dicotomica}
\KwInput{Funzione obiettivo $f$, oracolo di derivata $f'$, estremi iniziali $[x_-, x_+]$ con $f'(x_-) < 0$ e $f'(x_+) > 0$, tolleranza gradiente $\varepsilon > 0$, ampiezza minima $\delta > 0$}
\KwOutput{Punto di minimo locale approssimato $x^*$}

\While{$(x_+ - x_- > \delta)$}{
    $x \gets \frac{x_- + x_+}{2}$\;
    $g \gets f'(x)$ \tcp*{Valutazione oracolare della derivata}
    
    \If{$|g| \le \varepsilon$}{
        \Return{$x$} \tcp*{Trovato punto stazionario con $|f'(x)| \le \varepsilon$}
    }
    
    \eIf{$g < 0$}{
        $x_- \gets x$ \tcp*{La funzione scende a destra: scarta $[x_-, x]$}
    }{
        $x_+ \gets x$ \tcp*{La funzione sale a destra: scarta $[x, x_+]$}
    }
}
\Return{$\frac{x_- + x_+}{2}$}\;
\end{algorithm}

\subsection{Analisi della Convergenza Lineare e Complessit\`a Logaritmica}

\subsubsection{Riduzione Esatta dell'Intervallo}
Denotando con $D_0 = x_+ - x_-$ l'ampiezza iniziale dell'intervallo, a ogni iterazione $k$ l'ampiezza dell'intervallo di incertezza si dimezza esattamente:
\begin{equation}
    D_k = \frac{D_0}{2^k} = D_0 \cdot (0.5)^k.
\end{equation}
La sequenza delle ampiezze $\{D_k\}$ decade con un tasso di convergenza geometrico costante $r = 0.5$. Questo comportamento individua una \textbf{convergenza lineare} nello spazio di ricerca.

\subsubsection{Numero di Iterazioni in Funzione di $\delta$}
Affinch\'e l'intervallo si riduca al di sotto della soglia $\delta$:
\begin{equation}
    \frac{D_0}{2^k} \le \delta \iff 2^k \ge \frac{D_0}{\delta} \iff k \ge \log_2\left(\frac{D_0}{\delta}\right) = \frac{\log_{10}(D_0/\delta)}{\log_{10}(2)}.
\end{equation}
Poich\'e $\frac{1}{\log_{10}(2)} \approx 3.3219$:
\begin{equation}
    k \approx 3.32 \cdot \log_{10}\left(\frac{D_0}{\delta}\right).
    \label{eq:iterazioni-dicotomica-delta}
\end{equation}
Il significato computazionale della formula~\eqref{eq:iterazioni-dicotomica-delta} \`e limpido: \textbf{per guadagnare una cifra decimale esatta di risoluzione spaziale, la ricerca dicotomica richiede circa $3.32$ iterazioni}.

\subsubsection{Legame con l'Accuratezza del Gradiente $\varepsilon$ e $L$-Smoothness}
Supponiamo ora che la derivata prima $f'$ sia a sua volta Lipschitz-continua con costante $L_{f'}$ (ipotesi nota nella letteratura come \textbf{$L$-smoothness} di $f$):
\begin{equation}
    |f'(x) - f'(z)| \le L_{f'} |x - z|, \quad \forall x, z \in [x_-, x_+].
\end{equation}
Sia $x_*$ lo zero esatto di $f'$ ($f'(x_*) = 0$) contenuto nell'intervallo corrente di ampiezza $D_k \le \delta$. Poich\'e il punto medio $x$ dista dalla radice al pi\`u $\delta/2$:
\begin{equation}
    |f'(x)| = |f'(x) - f'(x_*)| \le L_{f'} |x - x_*| \le L_{f'} \frac{\delta}{2}.
\end{equation}
Per garantire che $|f'(x)| \le \varepsilon$, \`e dunque sufficiente che l'ampiezza dell'intervallo soddisfi:
\begin{equation}
    L_{f'} \frac{\delta}{2} \le \varepsilon \iff \delta = \frac{2\varepsilon}{L_{f'}}.
\end{equation}
Sostituendo tale valore di $\delta$ nella relazione di complessit\`a:
\begin{equation}
    k \ge \log_2\left( \frac{L_{f'} D_0}{2\varepsilon} \right) \approx 3.32 \cdot \log_{10}\left( \frac{L_{f'} D_0}{2\varepsilon} \right) \in \BigO{\log\left(\frac{L_{f'} D_0}{\varepsilon}\right)}.
    \label{eq:complessita-dicotomica-gradiente}
\end{equation}

\subsection{Confronto Epocale: Ricerca Globale vs.\ Ricerca Locale}
Mettendo a confronto diretto le formule di complessit\`a ricavate per i due approcci:

\begin{table}[H]
\centering
\small
\renewcommand{\arraystretch}{1.25}
\begin{tabularx}{\linewidth}{l Y Y}
\toprule
\textbf{Propriet\`a Computazionale} & \textbf{Ottimizzazione Globale (Griglia)} & \textbf{Ottimizzazione Locale (Dicotomica)} \\
\midrule
\textbf{Ipotesi di regolarit\`a} & $f$ \`e $L$-Lipschitziana su compatto & $f'$ \`e continua con cambio segno ($f \in C^1$) \\
\textbf{Garanzia teorica} & $\varepsilon$-minimo globale $f(\bar{x}) - f^* \le \varepsilon$ & Minimo locale con $|f'(x)| \le \varepsilon$ \\
\textbf{Complessit\`a asintotica} & $\BigO{\dfrac{L D}{\varepsilon}}$ (Sublineare) & $\BigO{\log\left(\dfrac{L_{f'} D}{\varepsilon}\right)}$ (Logaritmica / Lineare) \\
\midrule
\textbf{Valutazioni per $\varepsilon = 10^{-3}$} & $\approx 10^3$ campionamenti & $\approx 3.32 \times 3 \approx \mathbf{10\text{ iterazioni}}$ \\
\textbf{Valutazioni per $\varepsilon = 10^{-5}$} & $\approx 10^5$ campionamenti & $\approx 3.32 \times 5 \approx \mathbf{17\text{ iterazioni}}$ \\
\textbf{Valutazioni per $\varepsilon = 10^{-7}$} & $\approx 10^7$ (\textbf{10 Milioni!}) & $\approx 3.32 \times 7 \approx \mathbf{24\text{--}26\text{ iterazioni}}$ \\
\bottomrule
\end{tabularx}
\caption{Confronto di complessit\`a ed efficienza tra campionamento globale e ricerca dicotomica locale.}
\label{tab:confronto-globale-dicotomico}
\end{table}

\begin{noteblock}{Dimostrazione Sperimentale al Calcolatore}
Eseguendo lo script live di ricerca dicotomica (\texttt{dichotomic\_search}):
\begin{itemize}
    \item in appena \textbf{26 iterazioni}, l'algoritmo abbatte la norma della derivata al di sotto di $10^{-6}$ e localizza il punto di minimo con oltre 15 cifre decimali di precisione numerica;
    \item rispetto ai decine di milioni di campionamenti imposti dalla ricerca globale, si registra un'accelerazione computazionale nell'ordine di \textbf{$10^5\text{--}10^6$ volte}.
\end{itemize}
\`E l'effetto dirompente del passaggio da una complessit\`a algebrica inversa $\BigO{1/\varepsilon}$ a una complessit\`a logaritmica $\BigO{\log(1/\varepsilon)}$.
\end{noteblock}

\subsection{Inizializzazione dell'Intervallo: La Procedura di Bracketing}
La ricerca dicotomica richiede come pre-condizione fondamentale la disponibilit\`a di un intervallo iniziale $[x_-, x_+]$ che soddisfi la condizione di segno discorde $f'(x_-) < 0$ e $f'(x_+) > 0$. 

Se tale intervallo non \`e noto a priori, si ricorre a una procedura di \textbf{Bracketing} (espansione geometrica guidata):

\begin{algorithm}[H]
\caption{Procedura di Bracketing con Espansione Geometrica}
\label{alg:bracketing-espansione}
\KwInput{Punto di partenza $x_0$, oracolo di derivata $f'$, passo base $\Delta x > 0$}
\KwOutput{Estremo destro $x_+$ tale che $f'(x_+) > 0$}

$x_+ \gets x_0$\;
$\text{step} \gets \Delta x$\;
\While{$f'(x_+) \le 0$}{
    $x_+ \gets x_+ + \text{step}$\;
    $\text{step} \gets 2 \cdot \text{step}$ \tcp*{Raddoppio del passo (crescita geometrica)}
}
\Return{$x_+$}\;
\end{algorithm}
(Una procedura perfettamente simmetrica opera verso sinistra procedendo a passi negativi per determinare $x_-$ tale che $f'(x_-) < 0$).

\subsubsection{Casi Particolari e Patologie di Bracketing}
\begin{enumerate}
    \item \textbf{Funzioni coercive:} Se $f$ \`e una funzione coerciva ($\lim_{|x| \to \infty} f(x) = +\infty$), la derivata prima deve per forza assumere valori positivi per $x \to +\infty$ e negativi per $x \to -\infty$. La procedura di espansione garantisce l'arresto in un numero finito e modesto di raddoppi.
    \item \textbf{Funzioni illimitate inferiormente:} Se $f^* = -\infty$, la procedura diverger\`a a $\pm\infty$ senza mai arrestarsi; nei software di produzione \`e dunque obbligatorio impostare una soglia massima sul numero di raddoppi.
    \item \textbf{Controesempio patologico periodico:}
    Si consideri la funzione trigonometrica periodica:
    \begin{equation}
        f(x) = \sin\left(\pi x + \frac{3\pi}{4}\right) \implies f'(x) = \pi \cos\left(\pi x + \frac{3\pi}{4}\right).
    \end{equation}
    Inizializzando la procedura in $x_0 = 0$ con passo base $\Delta x = 1$ e raddoppiando a ogni iterazione, si generano punti di campionamento a distanze geometriche (o considerando le potenze pure $x_k = 2^k$ per $k \ge 1$):
    \begin{equation}
        x_k = 2^k, \quad k \in \N.
    \end{equation}
    Poich\'e per ogni $k \ge 1$ il valore $2^k$ \`e un numero intero pari, ponendo $2^k = 2m$ con $m \in \N$:
    \begin{equation}
        f'(2^k) = \pi \cos\left(2m \pi + \frac{3\pi}{4}\right) = \pi \cos\left(\frac{3\pi}{4}\right) = -\pi \frac{\sqrt{2}}{2} \approx -2.2214.
    \end{equation}
    La derivata valutata nei punti di campionamento risulta \textbf{costantemente negativa e invariata} a ogni passo ($\approx -2.22$). L'algoritmo non intercetta mai il cambio di segno e continua a raddoppiare all'infinito, nonostante la funzione sia limitata in $[-1, 1]$ e ricca di infiniti minimi locali.
\end{enumerate}

\FloatBarrier
\section{Sintesi e Prospettive: Possiamo Fare Meglio?}
\label{sec:univariata-chiusura-prospettive}

Il percorso teorico e algoritmico tracciato in questa lezione stabilisce i seguenti capisaldi:
\begin{enumerate}
    \item L'ottimizzazione globale pura su funzioni generiche \`e impossibile; con sole ipotesi di Lipschitzianit\`a su compatto, essa richiede una complessit\`a minimax sublineare intrattabile $\BigO{LD/\varepsilon}$.
    \item La transizione pragmatica all'ottimizzazione locale sfrutta la struttura unimodale dei bacini di attrazione e converte il problema nell'annullamento della derivata prima ($f'(x_*) = 0$).
    \item L'algoritmo di \textbf{Ricerca Dicotomica} dimezza l'intervallo di incertezza a ogni iterazione, garantendo la convergenza in sole $\BigO{\log(1/\varepsilon)}$ iterazioni (26 iterazioni per raggiungere una precisione di $10^{-7}$).
\end{enumerate}

La lezione si chiude con il quesito aperto formulato dal docente:
\begin{quote}
\textit{``Possiamo fare ancora meglio di 26 iterazioni? Possiamo localizzare il minimo con la stessa precisione in 12 o 18 iterazioni?''}
\end{quote}
La risposta \`e affermativa: sfruttando informazioni di ordine superiore (come la curvatura $f''$) o approssimazioni interpolanti polinomiali locali (interpolazione quadratica della secante, interpolazione cubica e il metodo di Newton univariato), sar\`a possibile accelerare ulteriormente il tasso di convergenza passando dalla convergenza lineare alla convergenza superlineare e quadratica.
